\section{Test Cases}
\label{chap:test_test_case}

This section presents a series of tests and their results which, in the developer's opinion, cover all aspects necessary to validate the operation of the project, both of the I2C controller and the LIS3DH.

\subsection{Test Cases}

\subsubsection{Test Case 1: Device Initialization and WHO\_AM\_I Verification}

\textbf{Objectives:} Verify that the LIS3DH device responds to I$^2$C transactions and 
returns the correct device ID.

\textbf{Procedure:}
\begin{enumerate}
  \item Initialize LIS3DH in Normal mode, $\pm$4g, 100~Hz.
  \item Read WHO\_AM\_I register (0x0F).
  \item \textbf{Expected:} \texttt{0x33} (LIS3DH device ID, datasheet §8.1).
\end{enumerate}

\textbf{Code:}
\begin{lstlisting}[language=C, caption={WHO\_AM\_I verification test case}, 
label={code:testcase1}]
static void testcase1() {
    printf("--- Test Case 1: WHO_AM_I Verification ---\n");
    lis3dh_config.fscale = LIS3DH_FS_4G; 
    lis3dh_config.mode = LIS3DH_MODE_NORMAL;
    // ... init ...
    HAL_StatusTypeDef status = lis3dh_init(&lis3dh_acc, &hi2c1, 
                                          LIS3DH_I2C_DEFAULT_ADDRESS, &lis3dh_config);
    if (status == HAL_OK) 
        printf("PASS: LIS3DH init OK, WHO_AM_I = 0x33\n");
    else 
        printf("FAIL: LIS3DH init failed\n");
}
\end{lstlisting}


%%%%%

\subsubsection{Test Case 2: Operating Mode Configuration}
\textbf{Objective:} Validate that CTRL\_REG1/CTRL\_REG4 correctly configure all 
operating modes.

\textbf{Procedure:} For each mode (Normal/High-Res/Low-Power): configure → read → verify.

\textbf{Code:}
\begin{lstlisting}[language=C, caption={Operating mode configuration test}, 
label={code:testcase2}]
static bool testcase21() {  // Normal mode
    lis3dh_set_mode(&lis3dh_acc, LIS3DH_MODE_NORMAL);
    lis3dh_mode_t mode = lis3dh_get_mode(&lis3dh_acc);
    return (mode == LIS3DH_MODE_NORMAL);
}
// Similar tests for HIGH_RES and LOW_POWER
\end{lstlisting}

%%%%%



\subsubsection{Test Case 3: Full-Scale Configuration}
\textbf{Objective:} Validate configuration of $\pm$2g/$\pm$4g/$\pm$8g/$\pm$16g ranges 
in CTRL\_REG4.

\textbf{Procedure:} For each full-scale range: configure → read → verify.

\textbf{Code:}
\begin{lstlisting}[language=C, caption={Full-scale range configuration test}, 
label={code:testcase3}]
static bool testcase31() {  
    lis3dh_set_full_scale(&lis3dh_acc, LIS3DH_FS_2G);
    lis3dh_fscale_t fs = lis3dh_get_full_scale(&lis3dh_acc);
    return (fs == LIS3DH_FS_2G);
}
// Similar tests for +-4g/+-8g/+-16g
\end{lstlisting}

%%%%%%%%%%%%%%%%%

\subsubsection{Test Case 4: Acceleration Data Injection and Conversion}
\textbf{Objective:} Verify complete data flow: QOM injection → registers → I$^2$C read 
→ firmware scaling.

\textbf{Test sequence:}
\begin{enumerate}
  \item Configure Normal mode, $\pm$4g (So = 8~mg/LSB).
  \item Inject via QEMU: \texttt{qom-set accel-x 1000}.
  \item Firmware reads via \texttt{lis3dh\_read\_acc()} → verifies 1000~mg.
\end{enumerate}

\textbf{Code:}
\begin{lstlisting}[language=C, caption={Acceleration data injection test}, 
label={code:testcase4}]
lis3dh_config.fscale = LIS3DH_FS_4G; 
lis3dh_config.mode = LIS3DH_MODE_NORMAL;
// ... config & enable DRDY ...
// QEMU: qom-set accel-x 1000
lis3dh_read_acc(&lis3dh_acc);
printf("INT! X: %.2f mg", lis3dh_acc.data->acc_mg[0]);  // Expected: 1000.00
\end{lstlisting}

%%%%%%%%%%%%%%%%%%%

\subsubsection{Test Case 5: Data-Ready Interrupt Generation}
\textbf{Objective:} Validate INT1 interrupt generation when new data is available.

\textbf{Configuration:} \texttt{lis3dh\_enable\_drdy\_interrupt()}.

\textbf{Test sequence:}
\begin{enumerate}
  \item Inject X/Y/Z data via QOM properties.
  \item Verify EXTI1 (PA1) triggers → \texttt{HAL\_GPIO\_EXTI\_Callback()}.
  \item Read data + INT1\_SRC (clears latched interrupt).
  \item INT1 returns low.
\end{enumerate}

\textbf{Code:}
\begin{lstlisting}[language=C, caption={Data-Ready interrupt handling}, 
label={code:testcase5}]
void HAL_GPIO_EXTI_Callback(uint16_t GPIO_Pin) {
    if (GPIO_Pin == LIS3DH_INT1_Pin) g_data_ready = true;
}
// Main loop
if (g_data_ready)
{
  g_data_ready = false;  // Clear flag

  // Read acceleration data
  status = lis3dh_read_acc(&lis3dh_acc);
  if (status == HAL_OK)
  {
    printf("INT! X: %.2f mg, Y: %.2f mg, Z: %.2f mg\n",
        lis3dh_acc.data->acc_mg[0],
        lis3dh_acc.data->acc_mg[1],
        lis3dh_acc.data->acc_mg[2]);
  }
  else
  {
    printf("STM32: Read acc ERROR after interrupt\n");
  }

  // Read INT1_SRC to verify/clear interrupt source
  uint8_t int_src = 0;
  lis3dh_read_int_source(&lis3dh_acc, &int_src);
  printf("INT1_SRC: 0x%02X\n", int_src);
}
\end{lstlisting}




























